\documentclass[11pt,a4paper]{article}
\usepackage{../common/d3report}
\usepackage{amsmath}
\usepackage{listings}
\usepackage{enumitem}
\usepackage{float}
\usepackage{array}

\lstset{
  basicstyle=\ttfamily\small,
  breaklines=true,
  frame=single,
  columns=fullflexible,
  keepspaces=true
}

\begin{document}
\makereporttitle
  {Concurrency and Seat-Allocation Report}
  {Deliverable 3 --- Reliability and Concurrency}
  {CS455 Campus Ride-Pooling System \quad | \quad Draft: experiment implementation committed; execution evidence pending}

\section{Objective and scope}

This report documents the concurrency strategy used for ride seat allocation and the experiment
implemented for the project's mandatory concurrency check (M17). The principal requirement is
NFR-RD-05 in the Rider non-functional requirements:

\begin{quote}
With at least 50 concurrent acceptance and leave operations on a ride with one available seat,
repeated 100 times: zero over-allocations, and the seat count always equals capacity minus
occupants.
\end{quote}

The corresponding Ride Owner requirements are NFR-RO-DEP-01 (strong consistency between ride
lifecycle and request decisions), NFR-RO-DEP-02 (concurrent joins and acceptances must not exceed
fixed capacity), NFR-RO-DEP-03 (atomic persistence of request decisions and ride state), and
NFR-RO-DEP-04 (locked rides reject normal membership changes).

The scope here is the implemented ride/request workflow in the Next.js application and PostgreSQL
database. It does not claim to test a deployed multi-instance service, a 50-connection database
pool, payment settlement, or physical vehicle allocation.

\section{Shared state and correctness invariant}

For a ride $r$, define:
\[
  O(r) = \left|\{s \in \texttt{riders} :
    s.\texttt{ride\_id}=r.\texttt{id} \land s.\texttt{left\_at IS NULL}\}\right|
\]
and let $C(r)$ be \texttt{rides.capacity\_snapshot}. The seat-allocation invariant is
\[
  1 \le O(r) \le C(r), \qquad
  \operatorname{available}(r)=C(r)-O(r).
\]
The owner is included among active occupants. Every other active occupant must have a corresponding
join request in the \texttt{accepted} state. A rider who leaves remains in the \texttt{riders}
history, but their non-null \texttt{left\_at} means that they no longer consume a seat.

The contended resource is the ride row and its active occupant set. Two acceptance transactions can
both read the same apparent free seat if the read/check/write sequence is not serialized. A
second hazard is an acceptance racing with a leave: both operations change the same occupancy
count, and their ordering determines how many seats are free at each point.

\section{Implemented concurrency control}

\subsection{Transaction and lock sequence}

The request-decision route begins a PostgreSQL transaction and calls
\texttt{actOnJoinRequest}. The current implementation follows this sequence:

\begin{enumerate}
  \item Read the request's rider identifier.
  \item Lock that student row with \texttt{SELECT ... FOR UPDATE}.
  \item Lock the parent ride row with \texttt{SELECT ... FOR NO KEY UPDATE}.
  \item Lock and re-read the join-request row with \texttt{FOR NO KEY UPDATE}.
  \item Re-check authorization, request state, ride state/lock time, and capacity using the
        current active-seat rows.
  \item Apply the seat/request/event/notification changes in the same transaction.
  \item Commit successful state transitions; roll back rejected transitions and exceptions.
\end{enumerate}

The critical capacity check in the acceptance branch occurs after the ride lock is acquired:
the application reads active seats and refuses acceptance when their count is already at
\texttt{capacity\_snapshot}. The successful path inserts or reactivates the rider row, changes
the request to \texttt{accepted}, records its event, and performs related notifications and
fare/itinerary refreshes before the route commits.

The leave branch takes the same user/ride/request locks. It marks the rider row as left, changes
the request to \texttt{left}, records the event, updates notifications and refreshes the affected
fare/itinerary state in the same transaction. As both operations lock the same ride row before
reading and changing occupancy, competing operations on one ride are serialized. A transaction
that obtains the lock later observes the committed occupancy change and revalidates against it.

\subsection{Database-level guard}

A normal PostgreSQL \texttt{CHECK} constraint cannot safely express an aggregate over rows in a
different table. The schema therefore now also installs a trigger on activation of a
\texttt{riders} row. It locks the parent ride row with \texttt{FOR NO KEY UPDATE}, counts active
occupants, and rejects an insert/reactivation that would exceed capacity using SQLSTATE
\texttt{23514} and constraint name \texttt{riders\_ride\_capacity}. A second trigger rejects a
capacity reduction below the current active-occupant count.

This guard protects direct SQL writes and future write paths as well as the current API. It
complements, rather than replaces, the transaction-level check. The API still performs its
business validation and returns the appropriate application error; the database trigger is the
last line of defence against an over-capacity write.

\subsection{What the schema does and does not guarantee}

The schema's \texttt{capacity\_snapshot > 0} check validates that the stored capacity is positive;
by itself it does not compare capacity with the number of active \texttt{riders}. The new trigger
provides that cross-row enforcement. The unique constraint on
\texttt{riders(ride\_id,user\_id)} prevents duplicate historical seat rows for the same rider and
ride; the \texttt{left\_at} column distinguishes active from departed occupants.

\section{Race scenarios and expected outcomes}

\begin{longtable}{p{0.24\linewidth}p{0.34\linewidth}p{0.34\linewidth}}
\toprule
\textbf{Scenario} & \textbf{Hazard without serialization} & \textbf{Expected result with current design}\\
\midrule
Two acceptances compete for the last seat &
Both transactions read one free seat and both insert an active occupant. &
The ride-row lock serializes the checks. One acceptance commits; the next sees a full ride and is rejected.\\
\addlinespace
Acceptance races with leave &
A stale occupancy read can cause an over-allocation or a lost seat update. &
Both actions serialize on the ride row. The later action re-reads occupancy after the earlier commit.\\
\addlinespace
Two leaves race on one request &
Duplicate leave events or inconsistent request/seat state. &
The request and ride locks serialize the actions; the first transition commits and the later action sees a non-accepted request.\\
\addlinespace
Acceptance races with ride lock &
A rider may be accepted after membership and fare shares have been frozen. &
Both operations use the ride lock. Whichever commits first determines whether the acceptance is valid; after lock, normal acceptance is refused.\\
\addlinespace
Direct SQL seat insert races with another insert &
Application-only validation can be bypassed by another writer. &
The database trigger takes the ride-row lock and checks the active count before allowing either insert.\\
\bottomrule
\end{longtable}

\section{Experiment methodology}

\subsection{Reproducible test}

The integration test is
\texttt{apps/web/test/integration/concurrency/last-seat.int.test.ts}. It invokes the real Next.js
route handlers using the repository's integration-test request harness and a disposable local
PostgreSQL test database. It does not mock the database or replace the request-decision logic.

Each measured round begins with an auto ride of capacity three: the owner and one accepted rider
occupy two seats, leaving exactly one available seat. The test dispatches 50 operations together:
49 distinct riders attempt acceptance and the current accepted rider attempts to leave. The test
runs 100 rounds. When the leave obtains the ride lock first, two acceptance attempts may use the
two seats that become free; one of those new occupants is then removed serially to restore the
one-free-seat starting condition for the next round. This serial cleanup is outside the 50
operations measured in that round.

The test uses a one-rupee fare and a \texttt{from\_hub} ride so the request snapshot does not change
materially as occupants leave or join. This isolates the capacity race from the separate
reconfirmation policy. Each round uses 49 distinct requester identities; pending requests are
reused where possible and recreated after a rider has left.

For each completed route operation, the test reads committed database state and checks:
\begin{itemize}
  \item active occupants never exceed the ride capacity;
  \item available seats equal capacity minus active occupants;
  \item every non-owner active occupant has exactly one accepted request;
  \item the round ends with two active occupants in a capacity-three ride (one available seat).
\end{itemize}
The test also includes a separate direct-database-writer race: two inserts compete for the final
seat and exactly one must be rejected by the database trigger.

\subsection{Concurrency interpretation and limits}

The 50 route-handler invocations are launched together with \texttt{Promise.allSettled}; the
integration harness invokes handlers directly rather than making network requests to a running
HTTP server. The application database pool is configured with \texttt{max: 10}. Consequently,
this experiment creates 50 concurrent application-level operations, but at most ten connections
from this pool can be active at once. It must not be described as a 50-database-session test.

The test checks committed snapshots after route calls resolve. Because the observation query
uses the same ten-connection pool, it can be scheduled after additional concurrent transactions
have committed; these snapshots are not a linear trace of every individual commit. The trigger
provides write-time enforcement independently of this sampling. The test is a repeatable
correctness stress test of the actual request handlers and database, not a latency benchmark or
a proof about all possible schedules.

The repository's pull-request CI also provisions PostgreSQL 16 and runs the complete integration
suite. The new test therefore participates in that database-backed CI job; its result should be
recorded from the actual check run rather than inferred from the source code.

\subsection{Run instructions}

Run from \texttt{apps/web} after configuring the disposable local PostgreSQL database exactly as
described in \texttt{apps/web/setup.md}. The integration harness deliberately refuses to reset a
database unless its name ends in \texttt{\_test}, it is local, and
\texttt{ALLOW\_TEST\_DATABASE\_RESET} matches that database name.

\begin{lstlisting}
npm run test:integration -- test/integration/concurrency/last-seat.int.test.ts
\end{lstlisting}

The experiment prints a JSON summary prefixed by
\texttt{[NFR-RD-05 concurrency results]}. Preserve that output with the test-run log when
recording evidence. The direct-writer test is included in the same file.

\section{Results}

\textbf{Execution status: passed locally on 4 November 2026.} The complete
PostgreSQL-backed integration suite passed: 16 test files and 112 tests, with zero failures, in
53.27 seconds. The dedicated NFR-RD-05 test also passed both of its tests in that run.

The emitted result for the full-suite run was:
\begin{lstlisting}
{
  "iterations": 100,
  "concurrentOperationsPerRound": 50,
  "concurrentAcceptanceAttempts": 4900,
  "concurrentLeaveOperations": 100,
  "totalConcurrentOperations": 5000,
  "acceptedInConcurrentRounds": 170,
  "serialCleanupLeaves": 70,
  "maximumObservedOccupants": 3,
  "capacity": 3,
  "invariantViolations": 0
}
\end{lstlisting}

Thus, 5,000 mixed route-handler operations were launched in 100 rounds. The ride never exceeded
its capacity of three occupants, all measured occupancy snapshots satisfied the seat/request
invariants, and the round-end state was restored to two occupants with one seat available. The
separate direct-database-writer test also passed: exactly one of two competing inserts for the
final seat succeeded, the other was rejected with SQLSTATE \texttt{23514}, and reducing capacity
below active occupancy was rejected.

A separate targeted run also passed the same two tests. Its race-dependent counters differed
(166 concurrent acceptances and 66 serial cleanup leaves), as expected from lock acquisition
order; the invariant results were unchanged.

\section{Concurrency strategy alternatives}

\begin{longtable}{p{0.20\linewidth}p{0.37\linewidth}p{0.34\linewidth}}
\toprule
\textbf{Strategy} & \textbf{Advantages} & \textbf{Trade-offs for this workflow}\\
\midrule
Pessimistic row locking (implemented) &
Simple reasoning for a small, highly contended resource; the capacity check and state transition
share one transaction; no retry protocol is required for ordinary contention. &
Requests for the same ride wait behind its row lock. A long transaction reduces throughput for
that ride, so transactions must remain short and must not perform slow external work while holding
the lock.\\
\addlinespace
Optimistic version check &
Avoids holding a lock while reading; can perform well when conflicts are rare. &
Requires a version column or conditional update, conflict detection, and bounded retries. A
last-seat race is a high-conflict workload, so many callers may retry and still need a final
atomic capacity check.\\
\addlinespace
Advisory locks &
Can serialize a logical ride identifier without relying on a particular row lock. &
The application must define and consistently acquire the lock key on every write path. Advisory
locks are easier to omit accidentally and are less visible in ordinary row-level data operations.\\
\addlinespace
Serializable isolation &
The database detects executions that cannot be serialized and aborts transactions that conflict. &
Serialization failures require retry handling. It may abort unrelated work and adds complexity
without removing the need to define transaction boundaries and handle retries correctly.\\
\addlinespace
Queue serialization &
A single worker or partitioned queue can process changes for one ride in order. &
Adds queue infrastructure, asynchronous completion semantics, operational monitoring, and
recovery/replay concerns. It is more infrastructure than this synchronous request workflow needs.\\
\bottomrule
\end{longtable}

Pessimistic row locking is appropriate here because the invariant is tied to one ride, the
contended operation is short, and the product needs an immediate accept/reject response. The
database trigger closes the bypass around the application transaction. If contention becomes a
measured throughput problem, optimistic retries or partitioned queue processing can be evaluated
against workload evidence rather than adopted pre-emptively.

\section{Limitations and follow-up}

\begin{enumerate}
  \item Retain the local integration-test output as execution evidence for the D3 submission.
  \item Re-run the integration suite after any subsequent changes to the ride/request transaction
        or the capacity triggers.
  \item The current pool limit is ten connections. A separate deployment/load test is required
        before making claims about 50 concurrent database sessions or production throughput.
  \item The test does not simulate process crashes, multiple application hosts, network latency,
        or a production reverse proxy.
  \item The report covers seat capacity for a ride. Physical vehicle double-booking is a separate
        invariant and is not demonstrated by this experiment.
\end{enumerate}

\section{Conclusion}

The application serializes acceptance and leave decisions on the parent ride row and revalidates
occupancy inside the transaction. The schema-level trigger extends that protection to direct
database writers. The 100-round local PostgreSQL experiment and the complete integration suite
passed, with zero observed capacity-invariant violations. This is correctness evidence for the
implemented route-handler and database path; it is not a 50-database-session load test or a claim
about production throughput.

\end{document}
